% Beamer tech-talk template — engineering talk (~30 min, ~20 frames + backup).
% Shape: hook · problem · constraints · prior art · design · implementation · demo · results · lessons · future.
% Compile: latexmk -pdf -interaction=nonstopmode -halt-on-error main.tex
\documentclass[aspectratio=169,12pt]{beamer}

\usepackage[utf8]{inputenc}
\usepackage[T1]{fontenc}
\usepackage[english]{babel}
\usepackage{graphicx}
\usepackage{booktabs}
\usepackage{listings}
\usepackage{xcolor}
\usepackage{hyperref}

\usetheme{Minimal}
\setbeamertemplate{navigation symbols}{}

\title[<SHORT TITLE>]{<TALK TITLE>}
\subtitle{<SUBTITLE: clarify scope or hook>}
\author[<SHORT-NAME>]{<NAME>}
\institute[<ORG>]{<ROLE> @ <ORG>\\<VENUE>}
\date{\today}

\begin{document}

\begin{frame}[plain]
  \titlepage
\end{frame}

\begin{frame}{The hook}
  \vfill
  \begin{quote}
    \Large <The war story, the surprising number, the "wait what" moment that frames the talk>
  \end{quote}
  \vfill
  \note{60s. Visual > bullets here.}
\end{frame}

\begin{frame}{The problem}
  \begin{itemize}
    \item <One-sentence problem statement>
    \item <Why it matters --- who feels the pain, how often, how bad>
    \item <Why it's hard --- the constraint that rules out the obvious solution>
  \end{itemize}
\end{frame}

\begin{frame}{Why this is hard}
  \begin{itemize}
    \item \textbf{<CONSTRAINT 1>} --- one-line explanation
    \item \textbf{<CONSTRAINT 2>} --- one-line explanation
    \item \textbf{<CONSTRAINT 3>} --- one-line explanation
  \end{itemize}
\end{frame}

\begin{frame}{Prior art}
  \begin{itemize}
    \item \textbf{<APPROACH A>} --- what it does well, what it doesn't
    \item \textbf{<APPROACH B>} --- what it does well, what it doesn't
    \item \textbf{<APPROACH C>} --- what it does well, what it doesn't
  \end{itemize}
\end{frame}

\begin{frame}{Our approach (high level)}
  \begin{columns}[T]
    \begin{column}{0.55\textwidth}
      \begin{itemize}
        \item <One-sentence headline>
        \item <Key design decision 1>
        \item <Key design decision 2>
      \end{itemize}
    \end{column}
    \begin{column}{0.4\textwidth}
      % \includegraphics[width=\textwidth]{figures/REPLACE-architecture}
      \fbox{\rule{0pt}{4.5cm}\rule{\textwidth}{0pt}}
      \begin{center}\footnotesize REPLACE-architecture\end{center}
    \end{column}
  \end{columns}
\end{frame}

\begin{frame}{Design deep-dive: <COMPONENT 1>}
  \begin{columns}[T]
    \begin{column}{0.5\textwidth}
      \begin{itemize}
        \item Inputs / outputs
        \item The interesting bit --- the algorithm or data structure
        \item Why it's correct or why it scales
      \end{itemize}
    \end{column}
    \begin{column}{0.45\textwidth}
      \fbox{\rule{0pt}{4.5cm}\rule{\textwidth}{0pt}}
      \begin{center}\footnotesize REPLACE-component-1\end{center}
    \end{column}
  \end{columns}
\end{frame}

\begin{frame}{Design deep-dive: <COMPONENT 2>}
  \begin{itemize}
    \item Inputs / outputs
    \item The trade-off we made and why
    \item The thing we discovered while building this
  \end{itemize}
\end{frame}

\begin{frame}{Implementation notes}
  \begin{itemize}
    \item \textbf{Language / runtime:} <choice and why>
    \item \textbf{Built vs reused:} <what we wrote, what we leveraged>
    \item \textbf{Operational surface:} <how it's deployed, scaled, observed>
  \end{itemize}
\end{frame}

\begin{frame}{Live demo}
  \begin{center}
    % \includegraphics[width=0.85\textwidth]{figures/REPLACE-demo}
    \fbox{\rule{0pt}{5cm}\rule{0.85\textwidth}{0pt}}\\
    \footnotesize Live demo or asciinema GIF (REPLACE-demo.gif)
  \end{center}
  \note{3-5 minutes. Second-most memorable thing after the hook.}
\end{frame}

\begin{frame}{Results: <HEADLINE METRIC>}
  \begin{columns}[T]
    \begin{column}{0.5\textwidth}
      \begin{block}{Finding}
        \textbf{<NUMBER> <UNIT>} --- vs.\ <BASELINE> at <FACTOR>$\times$
      \end{block}
      \begin{itemize}
        \item Test conditions: <briefly>
        \item Statistical significance: <one line>
      \end{itemize}
    \end{column}
    \begin{column}{0.45\textwidth}
      \fbox{\rule{0pt}{4.5cm}\rule{\textwidth}{0pt}}
      \begin{center}\footnotesize REPLACE-headline-metric\end{center}
    \end{column}
  \end{columns}
\end{frame}

\begin{frame}{Results: <SUPPORTING METRIC>}
  \begin{itemize}
    \item <Sub-metric 1>: <number>
    \item <Sub-metric 2>: <number>
    \item <Caveat / scope of validity>
  \end{itemize}
\end{frame}

\begin{frame}{What we got wrong}
  \begin{itemize}
    \item \textbf{<MISTAKE 1>} --- what we thought, what actually happened, what we changed
    \item \textbf{<MISTAKE 2>} --- same shape
  \end{itemize}
\end{frame}

\begin{frame}{Lessons learned}
  \begin{itemize}
    \item <LESSON 1 --- generalizable beyond this project>
    \item <LESSON 2 --- generalizable beyond this project>
    \item <LESSON 3 --- generalizable beyond this project>
  \end{itemize}
\end{frame}

\begin{frame}{Future directions}
  \begin{itemize}
    \item <DIRECTION 1>
    \item <DIRECTION 2>
    \item <DIRECTION 3>
  \end{itemize}
\end{frame}

\begin{frame}{Resources}
  \begin{itemize}
    \item \textbf{Code:} \url{<URL>}
    \item \textbf{Paper / writeup:} \url{<URL>}
    \item \textbf{Slides:} \url{<URL>}
    \item \textbf{Contact:} <EMAIL> $\cdot$ <SOCIAL HANDLE>
  \end{itemize}
\end{frame}

\begin{frame}[plain]{Questions?}
  \vfill
  \centering
  \Large <EMAIL>
  \vfill
\end{frame}

% ─── Backup section ─────────────────────────────────────────────────────────
\appendix

\section*{Backup}

\begin{frame}[plain]{Backup slides}
  \vfill
  \centering
  \Large For Q\&A --- alternative designs, benchmarks, FAQ
  \vfill
\end{frame}

\begin{frame}{Backup: alternative design considered}
  \begin{itemize}
    \item What we considered
    \item Why we didn't pick it
    \item Conditions under which it would win
  \end{itemize}
\end{frame}

\begin{frame}{Backup: detailed benchmarks}
  \begin{itemize}
    \item Workload A: <numbers>
    \item Workload B: <numbers>
    \item Workload C: <numbers>
  \end{itemize}
\end{frame}

\begin{frame}{Backup: FAQ}
  \begin{quote}
    <Anticipated question>
  \end{quote}
  <One-paragraph answer>
\end{frame}

\end{document}
